\chapter{Devoirs surveillés}
Deux versions sont proposées pour chaque chapitre. La structure de référence est : automatismes sur 6 points, exercice sur 6 points et problème sur 8 points.
\section*{Chapitre 1 -- Suites numériques}
\dsheader{1}{A}{Suites numériques}{modéliser, calculer, raisonner, communiquer}
\begin{exercice}[1 -- Automatismes]\points{6}
\begin{enumerate}
\item Donner le coefficient multiplicateur associé à une hausse de $5\,\%$.
\item La suite $12,17,22,27,\ldots$ est-elle arithmétique ? Donner sa raison.
\item Combien de termes contient $u_4+u_5+\cdots+u_{18}$ ?
\item Si $w_0=80$ et $w_{n+1}=0{,}9w_n$, calculer $w_2$.
\item Écrire $1+2+\cdots+20$ à l'aide du symbole $\sum$.
\end{enumerate}
\end{exercice}
\begin{exercice}[2 -- Abonnements]\points{6}
Une plateforme compte $120\,000$ abonnés au début de l'année et en gagne $11\,000$ chaque mois. On note $u_n$ le nombre d'abonnés, en milliers, après $n$ mois.
\begin{enumerate}
\item Justifier que $(u_n)$ est arithmétique et préciser sa raison.
\item Exprimer $u_n$ en fonction de $n$, puis calculer $u_{12}$.
\item Calculer $u_1+\cdots+u_{12}$. Interpréter cette somme avec prudence.
\end{enumerate}
\end{exercice}
\begin{exercice}[3 -- Placement et objectif]\points{8}
Un capital initial de $650\,000$ Ar augmente de $5\,\%$ par an. On note $C_n$ le capital après $n$ années.
\begin{enumerate}
\item Écrire $C_{n+1}$ en fonction de $C_n$, puis $C_n$ en fonction de $n$.
\item Calculer $C_5$, arrondi à l'ariary.
\item Déterminer à la calculatrice le premier entier $n$ pour lequel $C_n$ dépasse $1\,000\,000$ Ar.
\item Expliquer pourquoi le modèle est géométrique et non arithmétique.
\end{enumerate}
\end{exercice}
\enseignant{\textbf{Barème.} Ex. 1 : 6 points répartis sur les automatismes. Ex. 2 : 6 points. Ex. 3 : 8 points.\\
\textbf{Éléments de correction.} Le coefficient est $1{,}05$. La première suite flash est arithmétique de raison $5$ ; la somme contient $15$ termes ; $w_2=64$ et $\sum_{k=1}^{20}k$. Pour l'exercice 2, $u_n=120+11n$, donc $u_{12}=252$ ; la somme demandée utilise $12\dfrac{u_1+u_{12}}2$ et cumule des états mensuels, pas un nombre de personnes distinctes. Pour l'exercice 3, $C_n=650\,000\times\left(1+\dfrac{5}{100}\right)^n$ ; contrôler le seuil par deux valeurs consécutives.}

\dsheader{1}{B}{Suites numériques}{modéliser, calculer, raisonner, communiquer}
\begin{exercice}[1 -- Automatismes]\points{6}
\begin{enumerate}
\item Donner le coefficient multiplicateur associé à une hausse de $4\,\%$.
\item La suite $12,17,22,27,\ldots$ est-elle arithmétique ? Donner sa raison.
\item Combien de termes contient $u_4+u_5+\cdots+u_{18}$ ?
\item Si $w_0=80$ et $w_{n+1}=0{,}9w_n$, calculer $w_2$.
\item Écrire $1+2+\cdots+20$ à l'aide du symbole $\sum$.
\end{enumerate}
\end{exercice}
\begin{exercice}[2 -- Abonnements]\points{6}
Une plateforme compte $145\,000$ abonnés au début de l'année et en gagne $9\,000$ chaque mois. On note $u_n$ le nombre d'abonnés, en milliers, après $n$ mois.
\begin{enumerate}
\item Justifier que $(u_n)$ est arithmétique et préciser sa raison.
\item Exprimer $u_n$ en fonction de $n$, puis calculer $u_{12}$.
\item Calculer $u_1+\cdots+u_{12}$. Interpréter cette somme avec prudence.
\end{enumerate}
\end{exercice}
\begin{exercice}[3 -- Placement et objectif]\points{8}
Un capital initial de $720\,000$ Ar augmente de $4\,\%$ par an. On note $C_n$ le capital après $n$ années.
\begin{enumerate}
\item Écrire $C_{n+1}$ en fonction de $C_n$, puis $C_n$ en fonction de $n$.
\item Calculer $C_5$, arrondi à l'ariary.
\item Déterminer à la calculatrice le premier entier $n$ pour lequel $C_n$ dépasse $1\,000\,000$ Ar.
\item Expliquer pourquoi le modèle est géométrique et non arithmétique.
\end{enumerate}
\end{exercice}
\enseignant{\textbf{Barème.} Ex. 1 : 6 points répartis sur les automatismes. Ex. 2 : 6 points. Ex. 3 : 8 points.\\
\textbf{Éléments de correction.} Le coefficient est $1{,}04$. La première suite flash est arithmétique de raison $5$ ; la somme contient $15$ termes ; $w_2=64$ et $\sum_{k=1}^{20}k$. Pour l'exercice 2, $u_n=145+9n$, donc $u_{12}=253$ ; la somme demandée utilise $12\dfrac{u_1+u_{12}}2$ et cumule des états mensuels, pas un nombre de personnes distinctes. Pour l'exercice 3, $C_n=720\,000\times\left(1+\dfrac{4}{100}\right)^n$ ; contrôler le seuil par deux valeurs consécutives.}

\section*{Chapitre 2 -- Fonctions exponentielles}
\dsheader{2}{A}{Fonctions exponentielles}{calculer, représenter, modéliser, interpréter}
\begin{exercice}[1 -- Puissances et variations]\points{6}
\begin{enumerate}
\item Simplifier $2^4\times2^{-7}$ et $\dfrac{5^{x+3}}{5^x}$.
\item Écrire $3^{-2}$ sous forme de fraction.
\item Donner le sens de variation de $f(x)=1{,}04^x$ et de $g(x)=-2\times0{,}75^x$.
\item Expliquer pourquoi $a^x+a^y=a^{x+y}$ est fausse en général.
\end{enumerate}
\end{exercice}
\begin{exercice}[2 -- Fréquentation d'un site]\points{6}
Au lancement, un site reçoit $2400$ visites par jour. La fréquentation augmente de $4\,\%$ par semaine. On modélise par $V(t)=2400\times1{,}04^t$.
\begin{enumerate}
\item Interpréter les nombres $2400$ et $1{,}04$.
\item Calculer $V(8)$, arrondi à l'unité.
\item Déterminer le premier nombre entier de semaines pour lequel $V(t)>3100$.
\end{enumerate}
\end{exercice}
\begin{exercice}[3 -- Taux moyen]\points{8}
Le chiffre d'affaires d'une entreprise passe de $82\,000$ à $104000$ euros en quatre ans.
\begin{enumerate}
\item Calculer le coefficient global.
\item Écrire le coefficient annuel moyen sous la forme d'une puissance.
\item Calculer le taux annuel moyen, arrondi à $0{,}01\,\%$.
\item Expliquer pourquoi on ne divise pas simplement le taux global par quatre.
\end{enumerate}
\end{exercice}
\enseignant{\textbf{Barème :} $6+6+8$. \textbf{Correction :} $2^4\times2^{-7}=2^{-3}=\dfrac18$, $\dfrac{5^{x+3}}{5^x}=125$, $3^{-2}=\dfrac19$. $f$ est croissante et $g$ croissante. Dans le problème, calculer $V(8)$ puis tester deux entiers consécutifs autour du seuil. Pour le taux moyen, $q=\left(\dfrac{V_f}{V_i}\right)^{\frac14}$ et $t=q-1$ ; les coefficients se multiplient.}

\dsheader{2}{B}{Fonctions exponentielles}{calculer, représenter, modéliser, interpréter}
\begin{exercice}[1 -- Puissances et variations]\points{6}
\begin{enumerate}
\item Simplifier $2^4\times2^{-7}$ et $\dfrac{5^{x+3}}{5^x}$.
\item Écrire $3^{-2}$ sous forme de fraction.
\item Donner le sens de variation de $f(x)=1{,}03^x$ et de $g(x)=-2\times0{,}75^x$.
\item Expliquer pourquoi $a^x+a^y=a^{x+y}$ est fausse en général.
\end{enumerate}
\end{exercice}
\begin{exercice}[2 -- Fréquentation d'un site]\points{6}
Au lancement, un site reçoit $1850$ visites par jour. La fréquentation augmente de $3\,\%$ par semaine. On modélise par $V(t)=1850\times1{,}03^t$.
\begin{enumerate}
\item Interpréter les nombres $1850$ et $1{,}03$.
\item Calculer $V(8)$, arrondi à l'unité.
\item Déterminer le premier nombre entier de semaines pour lequel $V(t)>2250$.
\end{enumerate}
\end{exercice}
\begin{exercice}[3 -- Taux moyen]\points{8}
Le chiffre d'affaires d'une entreprise passe de $82\,000$ à $98000$ euros en quatre ans.
\begin{enumerate}
\item Calculer le coefficient global.
\item Écrire le coefficient annuel moyen sous la forme d'une puissance.
\item Calculer le taux annuel moyen, arrondi à $0{,}01\,\%$.
\item Expliquer pourquoi on ne divise pas simplement le taux global par quatre.
\end{enumerate}
\end{exercice}
\enseignant{\textbf{Barème :} $6+6+8$. \textbf{Correction :} $2^4\times2^{-7}=2^{-3}=\dfrac18$, $\dfrac{5^{x+3}}{5^x}=125$, $3^{-2}=\dfrac19$. $f$ est croissante et $g$ croissante. Dans le problème, calculer $V(8)$ puis tester deux entiers consécutifs autour du seuil. Pour le taux moyen, $q=\left(\dfrac{V_f}{V_i}\right)^{\frac14}$ et $t=q-1$ ; les coefficients se multiplient.}

\section*{Chapitre 3 -- Fonction logarithme décimal}
\dsheader{3}{A}{Logarithme décimal}{calculer, résoudre, raisonner, communiquer}
\begin{exercice}[1 -- Connaissances et calculs]\points{6}
\begin{enumerate}
\item Calculer $\log(1000)$, $\log(0{,}01)$ et $\log(10^{-4})$.
\item Pour $x,y>0$, développer $\log\!\left(\dfrac{100x^3}{y} \right)$.
\item Résoudre $10^t=28000$ et donner $t$ au centième.
\end{enumerate}
\end{exercice}
\begin{exercice}[2 -- Objectif de croissance]\points{6}
Une valeur est multipliée chaque année par $1{,}05$. On cherche le temps nécessaire pour qu'elle double.
\begin{enumerate}
\item Traduire le problème par une équation.
\item Résoudre cette équation à l'aide du logarithme décimal.
\item Donner le premier nombre entier d'années permettant d'atteindre ou dépasser le double.
\item Vérifier avec les deux puissances consécutives utiles.
\end{enumerate}
\end{exercice}
\begin{exercice}[3 -- Niveau sonore]\points{8}
Le niveau sonore est $L=10\log\!\left(\dfrac I{I_0}\right)$, en décibels.
\begin{enumerate}
\item Calculer $L$ lorsque $I=10^6I_0$.
\item Pour $L=84$, exprimer puis calculer le rapport $\dfrac I{I_0}$.
\item Comparer les intensités correspondant à $84$ dB et à $74$ dB.
\end{enumerate}
\end{exercice}
\enseignant{\textbf{Barème :} $6+6+8$. \textbf{Correction :} Les valeurs usuelles sont $3,-2,-4$. Le développement donne $2+3\log x-\log y$. $t=\log(28000)$. Pour le doublement, $x=\dfrac{\log2}{\log(1{,}05)}$ puis on prend l'entier supérieur et on contrôle. $\dfrac I{I_0}=10^{\frac L{10}}$ ; un écart de $10$ dB correspond à une intensité multipliée par $10$.}

\dsheader{3}{B}{Logarithme décimal}{calculer, résoudre, raisonner, communiquer}
\begin{exercice}[1 -- Connaissances et calculs]\points{6}
\begin{enumerate}
\item Calculer $\log(1000)$, $\log(0{,}01)$ et $\log(10^{-4})$.
\item Pour $x,y>0$, développer $\log\!\left(\dfrac{100x^3}{y} \right)$.
\item Résoudre $10^t=32000$ et donner $t$ au centième.
\end{enumerate}
\end{exercice}
\begin{exercice}[2 -- Objectif de croissance]\points{6}
Une valeur est multipliée chaque année par $1{,}07$. On cherche le temps nécessaire pour qu'elle double.
\begin{enumerate}
\item Traduire le problème par une équation.
\item Résoudre cette équation à l'aide du logarithme décimal.
\item Donner le premier nombre entier d'années permettant d'atteindre ou dépasser le double.
\item Vérifier avec les deux puissances consécutives utiles.
\end{enumerate}
\end{exercice}
\begin{exercice}[3 -- Niveau sonore]\points{8}
Le niveau sonore est $L=10\log\!\left(\dfrac I{I_0}\right)$, en décibels.
\begin{enumerate}
\item Calculer $L$ lorsque $I=10^6I_0$.
\item Pour $L=78$, exprimer puis calculer le rapport $\dfrac I{I_0}$.
\item Comparer les intensités correspondant à $78$ dB et à $68$ dB.
\end{enumerate}
\end{exercice}
\enseignant{\textbf{Barème :} $6+6+8$. \textbf{Correction :} Les valeurs usuelles sont $3,-2,-4$. Le développement donne $2+3\log x-\log y$. $t=\log(32000)$. Pour le doublement, $x=\dfrac{\log2}{\log(1{,}07)}$ puis on prend l'entier supérieur et on contrôle. $\dfrac I{I_0}=10^{\frac L{10}}$ ; un écart de $10$ dB correspond à une intensité multipliée par $10$.}

\section*{Chapitre 4 -- Fonction inverse}
\dsheader{4}{A}{Fonction inverse}{calculer, raisonner, représenter, communiquer}
\begin{exercice}[1 -- Fonction inverse]\points{6}
On considère $f(x)=2+\dfrac{16}x$.
\begin{enumerate}
\item Donner son domaine.
\item Calculer sa dérivée et étudier son signe.
\item Donner ses asymptotes et son sens de variation sur chaque intervalle du domaine.
\end{enumerate}
\end{exercice}
\begin{exercice}[2 -- Tableau de signes]\points{6}
Résoudre $\dfrac{(x-4)(x+2)}x\geq0$. Le tableau doit faire apparaître les zéros et la valeur interdite.
\end{exercice}
\begin{exercice}[3 -- Coût moyen]\points{8}
Pour $x>0$, un coût moyen est modélisé par $C(x)=x+\dfrac{16}x+12$.
\begin{enumerate}
\item Calculer $C'(x)$ et factoriser son numérateur.
\item Dresser le tableau de variations de $C$ sur $]0,+\infty[$.
\item Déterminer la valeur de $x$ qui minimise le coût moyen et calculer ce minimum.
\item Faire une esquisse cohérente de la courbe.
\end{enumerate}
\end{exercice}
\enseignant{\textbf{Barème :} $6+6+8$. \textbf{Correction :} $D_f=\R\setminus\{0\}$, $f'(x)=-\dfrac{16}{x^2}<0$, asymptotes $x=0$ et $y=2$. Pour le quotient, ordonner $-2,0,4$ et exclure $0$. Pour le coût, $C'(x)=1-\dfrac{16}{x^2}=\dfrac{(x-4)(x+4)}{x^2}$ ; minimum en $x=4$, de valeur $20$.}

\dsheader{4}{B}{Fonction inverse}{calculer, raisonner, représenter, communiquer}
\begin{exercice}[1 -- Fonction inverse]\points{6}
On considère $f(x)=2+\dfrac{9}x$.
\begin{enumerate}
\item Donner son domaine.
\item Calculer sa dérivée et étudier son signe.
\item Donner ses asymptotes et son sens de variation sur chaque intervalle du domaine.
\end{enumerate}
\end{exercice}
\begin{exercice}[2 -- Tableau de signes]\points{6}
Résoudre $\dfrac{(x-3)(x+2)}x\geq0$. Le tableau doit faire apparaître les zéros et la valeur interdite.
\end{exercice}
\begin{exercice}[3 -- Coût moyen]\points{8}
Pour $x>0$, un coût moyen est modélisé par $C(x)=x+\dfrac{9}x+12$.
\begin{enumerate}
\item Calculer $C'(x)$ et factoriser son numérateur.
\item Dresser le tableau de variations de $C$ sur $]0,+\infty[$.
\item Déterminer la valeur de $x$ qui minimise le coût moyen et calculer ce minimum.
\item Faire une esquisse cohérente de la courbe.
\end{enumerate}
\end{exercice}
\enseignant{\textbf{Barème :} $6+6+8$. \textbf{Correction :} $D_f=\R\setminus\{0\}$, $f'(x)=-\dfrac{9}{x^2}<0$, asymptotes $x=0$ et $y=2$. Pour le quotient, ordonner $-2,0,3$ et exclure $0$. Pour le coût, $C'(x)=1-\dfrac{9}{x^2}=\dfrac{(x-3)(x+3)}{x^2}$ ; minimum en $x=3$, de valeur $18$.}

\section*{Chapitre 5 -- Statistiques à deux variables quantitatives}
\dsheader{5}{A}{Statistiques à deux variables}{représenter, modéliser, calculer, interpréter}
\begin{exercice}[1 -- Construire et décrire]\points{6}
On observe les couples $(1;8)$, $(2;11)$, $(3;13)$, $(4;17)$, $(5;19)$ et $(6;23)$.
\begin{enumerate}
\item Représenter le nuage dans un repère correctement légendé.
\item Calculer les coordonnées du point moyen.
\item Décrire la tendance et la dispersion sans conclure à une causalité.
\end{enumerate}
\end{exercice}
\begin{exercice}[2 -- Exploiter un ajustement]\points{6}
Une droite d'ajustement est $y=3{,}2x+12$, établie pour $x\in[2,15]$.
\begin{enumerate}
\item Estimer $y$ pour $x=10$ et qualifier l'estimation.
\item Estimer $x$ lorsque $y=60$.
\item Estimer $y$ pour $x=40$, puis expliquer pourquoi ce résultat doit être considéré avec prudence.
\end{enumerate}
\end{exercice}
\begin{exercice}[3 -- Changement de variable]\points{8}
On pose $z=\log(y)$ et l'ajustement de $(x,z)$ est $z=0{,}06x+1{,}1$.
\begin{enumerate}
\item Exprimer $y$ en fonction de $x$.
\item Estimer $y$ pour $x=8$.
\item Expliquer l'intérêt du changement de variable et une limite de l'extrapolation.
\end{enumerate}
\end{exercice}
\enseignant{\textbf{Barème :} $6+6+8$. \textbf{Correction :} Calculer $\bar x$ et $\bar y$ à partir des six valeurs. Avec la droite, substituer $x$ ou résoudre une équation ; $x=10$ est une interpolation et $x=40$ une extrapolation éloignée. Le dernier modèle est $y=10^{0{,}06x+1{,}1}$.}

\dsheader{5}{B}{Statistiques à deux variables}{représenter, modéliser, calculer, interpréter}
\begin{exercice}[1 -- Construire et décrire]\points{6}
On observe les couples $(1;8)$, $(2;11)$, $(3;13)$, $(4;17)$, $(5;22)$ et $(6;23)$.
\begin{enumerate}
\item Représenter le nuage dans un repère correctement légendé.
\item Calculer les coordonnées du point moyen.
\item Décrire la tendance et la dispersion sans conclure à une causalité.
\end{enumerate}
\end{exercice}
\begin{exercice}[2 -- Exploiter un ajustement]\points{6}
Une droite d'ajustement est $y=2{,}4x+18$, établie pour $x\in[2,15]$.
\begin{enumerate}
\item Estimer $y$ pour $x=10$ et qualifier l'estimation.
\item Estimer $x$ lorsque $y=60$.
\item Estimer $y$ pour $x=40$, puis expliquer pourquoi ce résultat doit être considéré avec prudence.
\end{enumerate}
\end{exercice}
\begin{exercice}[3 -- Changement de variable]\points{8}
On pose $z=\log(y)$ et l'ajustement de $(x,z)$ est $z=0{,}06x+1{,}1$.
\begin{enumerate}
\item Exprimer $y$ en fonction de $x$.
\item Estimer $y$ pour $x=8$.
\item Expliquer l'intérêt du changement de variable et une limite de l'extrapolation.
\end{enumerate}
\end{exercice}
\enseignant{\textbf{Barème :} $6+6+8$. \textbf{Correction :} Calculer $\bar x$ et $\bar y$ à partir des six valeurs. Avec la droite, substituer $x$ ou résoudre une équation ; $x=10$ est une interpolation et $x=40$ une extrapolation éloignée. Le dernier modèle est $y=10^{0{,}06x+1{,}1}$.}

\section*{Chapitre 6 -- Probabilités conditionnelles}
\dsheader{6}{A}{Probabilités conditionnelles}{modéliser, calculer, raisonner, interpréter}
\begin{exercice}[1 -- Arbre pondéré]\points{6}
Une proportion 0{,}35 des commandes est express. Le taux de retard est 0{,}05 parmi les commandes express et 0{,}12 parmi les autres.
\begin{enumerate}
\item Définir les événements $E$ et $R$, puis construire l'arbre pondéré complet.
\item Calculer $P(E\cap R)$.
\item Calculer $P(R)$ par la formule des probabilités totales.
\item Calculer $P_R(E)$ et interpréter.
\end{enumerate}
\end{exercice}
\begin{exercice}[2 -- Indépendance]\points{6}
Dans une autre étude, $P(A)=0{,}4$, $P(B)=0{,}25$ et $P(A\cap B)=0{,}10$.
\begin{enumerate}
\item Calculer $P(A)P(B)$.
\item Les événements sont-ils indépendants ? Justifier.
\item Calculer $P_A(B)$ et confirmer la conclusion.
\end{enumerate}
\end{exercice}
\begin{exercice}[3 -- Tableau d'effectifs]\points{8}
Sur $800$ clients, $320$ ont une carte premium ; $96$ clients premium et $72$ clients non premium ont résilié. Construire le tableau complet puis calculer la probabilité d'être premium sachant que le client a résilié.
\end{exercice}
\enseignant{\textbf{Barème :} $8+6+6$. \textbf{Correction :} $P(R)=P(E)P_E(R)+P(\overline E)P_{\overline E}(R)$ puis $P_R(E)=\dfrac{P(E\cap R)}{P(R)}$. Pour l'indépendance, comparer l'intersection au produit $0{,}10$. Dans le tableau, il y a $168$ résiliations, donc $P_R(P)=\dfrac{96}{168}=\dfrac47$.}

\dsheader{6}{B}{Probabilités conditionnelles}{modéliser, calculer, raisonner, interpréter}
\begin{exercice}[1 -- Arbre pondéré]\points{6}
Une proportion 0{,}45 des commandes est express. Le taux de retard est 0{,}08 parmi les commandes express et 0{,}15 parmi les autres.
\begin{enumerate}
\item Définir les événements $E$ et $R$, puis construire l'arbre pondéré complet.
\item Calculer $P(E\cap R)$.
\item Calculer $P(R)$ par la formule des probabilités totales.
\item Calculer $P_R(E)$ et interpréter.
\end{enumerate}
\end{exercice}
\begin{exercice}[2 -- Indépendance]\points{6}
Dans une autre étude, $P(A)=0{,}4$, $P(B)=0{,}25$ et $P(A\cap B)=0{,}14$.
\begin{enumerate}
\item Calculer $P(A)P(B)$.
\item Les événements sont-ils indépendants ? Justifier.
\item Calculer $P_A(B)$ et confirmer la conclusion.
\end{enumerate}
\end{exercice}
\begin{exercice}[3 -- Tableau d'effectifs]\points{8}
Sur $800$ clients, $320$ ont une carte premium ; $96$ clients premium et $72$ clients non premium ont résilié. Construire le tableau complet puis calculer la probabilité d'être premium sachant que le client a résilié.
\end{exercice}
\enseignant{\textbf{Barème :} $8+6+6$. \textbf{Correction :} $P(R)=P(E)P_E(R)+P(\overline E)P_{\overline E}(R)$ puis $P_R(E)=\dfrac{P(E\cap R)}{P(R)}$. Pour l'indépendance, comparer l'intersection au produit $0{,}10$. Dans le tableau, il y a $168$ résiliations, donc $P_R(P)=\dfrac{96}{168}=\dfrac47$.}

\section*{Chapitre 7 -- Variables aléatoires discrètes et loi binomiale}
\dsheader{7}{A}{Variables aléatoires et loi binomiale}{modéliser, calculer, raisonner, interpréter}
\begin{exercice}[1 -- Espérance d'un jeu]\points{6}
Un jeu coûte $3$ euros. Le gain brut vaut $0$, $5$ ou $15$ euros avec probabilités $0{,}65$, $0{,}25$ et $0{,}10$.
\begin{enumerate}
\item Construire la loi du gain net $G$.
\item Calculer $E(G)$.
\item Le jeu est-il favorable au joueur ? Interpréter sans affirmer ce qui arrivera lors d'une partie.
\end{enumerate}
\end{exercice}
\begin{exercice}[2 -- Loi binomiale]\points{6}
Un client répond favorablement avec probabilité $0{,}30$. Un conseiller contacte indépendamment $10$ clients. On note $X$ le nombre de réponses favorables.
\begin{enumerate}
\item Justifier la loi de $X$ et donner ses paramètres.
\item Calculer $P(X=2)$.
\item Calculer $P(X=0)$ puis $P(X\geq1)$.
\item Calculer et interpréter $E(X)$.
\end{enumerate}
\end{exercice}
\begin{exercice}[3 -- Triangle de Pascal]\points{8}
\begin{enumerate}
\item Écrire la ligne $n=6$ du triangle de Pascal.
\item Lire $\binom62$ et $\binom64$.
\item Expliquer le rôle de $\binom nk$ dans la formule binomiale.
\end{enumerate}
\end{exercice}
\enseignant{\textbf{Barème :} $6+8+6$. \textbf{Correction :} Le gain net vaut $-3,2,12$ ; $E(G)=-1{,}95+0{,}50+1{,}20=-0{,}25$ euro. $X\sim\mathcal B(10,0{,}30)$ ; appliquer la formule, utiliser $1-P(X=0)$ pour « au moins un », et $E(X)=np$. La ligne $n=6$ est $1,6,15,20,15,6,1$.}

\dsheader{7}{B}{Variables aléatoires et loi binomiale}{modéliser, calculer, raisonner, interpréter}
\begin{exercice}[1 -- Espérance d'un jeu]\points{6}
Un jeu coûte $3$ euros. Le gain brut vaut $0$, $5$ ou $15$ euros avec probabilités $0{,}65$, $0{,}25$ et $0{,}10$.
\begin{enumerate}
\item Construire la loi du gain net $G$.
\item Calculer $E(G)$.
\item Le jeu est-il favorable au joueur ? Interpréter sans affirmer ce qui arrivera lors d'une partie.
\end{enumerate}
\end{exercice}
\begin{exercice}[2 -- Loi binomiale]\points{6}
Un client répond favorablement avec probabilité $0{,}25$. Un conseiller contacte indépendamment $8$ clients. On note $X$ le nombre de réponses favorables.
\begin{enumerate}
\item Justifier la loi de $X$ et donner ses paramètres.
\item Calculer $P(X=2)$.
\item Calculer $P(X=0)$ puis $P(X\geq1)$.
\item Calculer et interpréter $E(X)$.
\end{enumerate}
\end{exercice}
\begin{exercice}[3 -- Triangle de Pascal]\points{8}
\begin{enumerate}
\item Écrire la ligne $n=6$ du triangle de Pascal.
\item Lire $\binom62$ et $\binom64$.
\item Expliquer le rôle de $\binom nk$ dans la formule binomiale.
\end{enumerate}
\end{exercice}
\enseignant{\textbf{Barème :} $6+8+6$. \textbf{Correction :} Le gain net vaut $-3,2,12$ ; $E(G)=-1{,}95+0{,}50+1{,}20=-0{,}25$ euro. $X\sim\mathcal B(8,0{,}25)$ ; appliquer la formule, utiliser $1-P(X=0)$ pour « au moins un », et $E(X)=np$. La ligne $n=6$ est $1,6,15,20,15,6,1$.}
